\section{Discussion and Conclusion}
\label{sec:discussion_conclusion}

\subsection{Discussion}
\label{sec:discussion}

This work evaluates frozen V-JEPA~2.1 image and video representations as a source of dense geometric and semantic information for aerial 3D scene understanding.  The experiments are intentionally diagnostic rather than purely architectural: the formulation in Eq.~\eqref{eq:task_functions} separates monocular depth estimation, pose-aware multiview depth estimation, and semantic scene completion, while Fig.~\ref{fig:evaluation_overview} and Table~\ref{tab:method_depth_variants} organize the study as a sequence of increasingly structured probes.  This design makes it possible to ask which components actually contribute to aerial metric depth and occupancy prediction: the frozen representation, the dense decoder, altitude/intrinsics conditioning, temporal aggregation, or explicit multiview correspondence.

\paragraph{Frozen V-JEPA~2.1 features contain useful aerial depth information, but dense structure and metadata are necessary.}
The linear-probe results in Table~\ref{tab:mde_comparison} show that some metric-depth information is directly accessible from the frozen V-JEPA~2.1 representation.  The RGB linear probe reaches an RMSE of 7.808~m, and the four-frame video linear probe improves this to 7.226~m.  This supports the view that the final V-JEPA~2.1 features encode nontrivial geometry even without backbone finetuning or a nonlinear spatial decoder.  However, the gap between the linear probes and the structured decoders is large.  SFP\_FiLM\_DPT reduces RMSE to 5.160~m, and FiLM\_DPT\_RGB gives the best overall V-JEPA~2.1 depth result with 4.621~m RMSE and 0.129 AbsRel.  Thus, the main conclusion from the depth study is not that depth is already linearly solved by V-JEPA~2.1, but that frozen features become substantially more useful when they are decoded with a dense prediction head and conditioned on aerial metadata.

The strongest evidence for metadata conditioning comes from the comparison between the plain DPT and FiLM-conditioned variants.  In Table~\ref{tab:mde_comparison}, the unconditioned DPT obtains 11.471~m RMSE, while FiLM\_DPT\_RGB obtains 4.621~m.  This difference should not be interpreted as a general failure of DPT-style decoding; rather, under this frozen-backbone aerial setting, the tested DPT path is insufficiently constrained unless the decoder receives scale-relevant metadata.  The FiLM mechanism in Eq.~\eqref{eq:metadata_vector} and Eq.~\eqref{eq:film} is therefore important because altitude and intrinsics directly affect the mapping from image evidence to metric depth.  A UAV observing the same object from 30~m, 40~m, and 50~m changes the apparent scale of the scene, and the frozen foundation representation is not explicitly trained to resolve this scale in OccuFly coordinates.  Conditioning the decoder with $h_t$ and $K_t$ gives the trainable head a physically meaningful signal for metric calibration.

\paragraph{Image and video features are complementary, but the current video pathway does not dominate overall.}
The video experiments give a more nuanced result.  The four-frame video linear probe improves over the RGB linear probe in Table~\ref{tab:mde_comparison}, suggesting that the frozen video tokens contain depth-relevant temporal information.  At 30~m altitude, Table~\ref{tab:mde_comparison_by_altitude} shows that DTA\_FiLM\_DPT\_Video achieves the best RMSE among the listed methods, with 2.939~m, outperforming FiLM\_DPT\_RGB at 3.448~m and DepthAnything2-OccuFly at 3.108~m.  This may indicate that local temporal aggregation can help when the target geometry is close enough for motion parallax and local frame-to-frame evidence to be informative.

The same trend does not hold uniformly.  At 40~m and 50~m, FiLM\_DPT\_RGB is stronger than DTA\_FiLM\_DPT\_Video among the V-JEPA~2.1 variants, and Table~\ref{tab:mde_comparison_by_scene} shows a similar scene-dependent pattern: the video FiLM variant is best on scene\_08, whereas the RGB FiLM variant is better on scene\_09.  This need not mean that video features are unhelpful, but that the current Dense Temporal Attention adapter in Eq.~\eqref{eq:dta} does not yet provide a consistently better target-frame depth representation than the RGB FiLM decoder.  Temporal tokens appear useful, but their benefit depends on altitude, scene content, camera motion, and the effectiveness of the adapter used to convert tubelet features into target-frame dense maps.

\paragraph{The altitude-interpolation result supports the use of physically meaningful conditioning, with a split caveat.}
Table~\ref{tab:occufly-depth-30-50} reports the altitude-interpolation result for V-JEPA~2.1 RGB FiLM\_DPT\_RGB.  The model reaches 0.9653 $\delta_1$, 0.1021 AbsRel, and 4.7071~m RMSE when evaluated at the held-out middle altitude.  This is consistent with the intended role of FiLM: the metadata pathway should not only memorize observed flight heights, but also help interpolate scale behavior when the test altitude lies between training altitudes. The result is strong evidence that altitude/intrinsics conditioning is promising.

\paragraph{Depth quality transfers to SSC only when the lifting architecture uses it effectively.}
The SSC results in Table~\ref{tab:ssc_occufly_comparison} show that better depth alone is not sufficient to guarantee better semantic scene completion.  On the overall ``All'' group, CGFormer with the DAV2 prior obtains 36.77 SC IoU and 3.05 SSC mIoU, while CGFormer with the V-JEPA~2.1 FiLM\_DPT\_RGB prior obtains a similar SC IoU of 36.51 but a higher SSC mIoU of 3.31.  This suggests that the V-JEPA-derived depth prior does not simply increase occupied/free-space completion, but can improve semantic occupancy when inserted into the CGFormer-style evaluation path.  The improvement is modest, but it is meaningful because it is obtained by replacing the depth prior rather than redesigning the full SSC system.

The strongest overall SSC result is obtained by FSSC\_RGB\_context\_DPT\_depth\_Film\_DPT, which reaches 37.68 SC IoU and 3.46 SSC mIoU in the ``All'' group. The result also shows that the best semantic completion setting is not necessarily the one with the most heavily conditioned context branch. The variant FSSC\_RGB\_context\_FiLM\_DPT\_depth\_FiLM\_DPT reaches only 35.47 SC IoU and 1.73 SSC mIoU overall, despite using FiLM-conditioned components.  This indicates that adding conditioning everywhere can disturb semantic-context features or create a mismatch between the context and geometry pathways.  In this setting, the more effective configuration is to use RGB context with a DPT-style context pathway and a FiLM-conditioned depth pathway, rather than to FiLM-modulate both branches indiscriminately.

\paragraph{Pose-aware MVS is promising, but the current implementation is sensitive to conditioning.}
The Pose-MVS branch is motivated by Eq.~\eqref{eq:mvs_projection}, Eq.~\eqref{eq:mvs_cost_volume}, and Eq.~\eqref{eq:mvs_depth_estimate}: if calibrated source views are available, projected feature matching across depth hypotheses should provide a stronger geometric prior than monocular appearance alone.  The SSC results show that this hypothesis is only partially realized.  The unconditioned video-context MVS variant, FSSC\_video\_context\_DTA\_DPT\_depth\_DTA\_mvs, performs poorly overall, with 24.55 SC IoU and 0.17 SSC mIoU.  By contrast, adding FiLM to the MVS depth logits improves the same video/MVS family to 32.55 SC IoU and 2.36 SSC mIoU.  This large recovery indicates that altitude/intrinsics information is also important for the multiview branch, not only for monocular depth.

At the same time, the FiLM-conditioned MVS variant remains below the best RGB FoundationSSC variant overall.  This suggests that pose-aware matching is not yet fully exploited.  Possible reasons include imperfect target/source pairing and deviation in the orientation of the picture taken in the considered sequence of frames.

\paragraph{SC IoU and SSC mIoU reveal different failure modes.}
Table~\ref{tab:ssc_occufly_comparison} also shows that occupied-space completion and semantic completion should be interpreted separately.  Several variants obtain competitive SC IoU while still having low SSC mIoU.  This means that the model may locate occupied regions reasonably well but still assign incorrect semantic labels.  This distinction is central for aerial scene understanding: collision avoidance and free-space estimation depend strongly on occupancy, while landing-zone assessment, human-UAV interaction, and object-aware planning require semantic correctness.  The per-class bar chart in Fig.~\ref{fig:occufly_ssc_per_class_iou_bar_chart} and the heatmap in Fig.~\ref{fig:occufly_ssc_per_class_iou_heatmap} should therefore be used alongside the aggregate numbers to identify which semantic categories dominate the mIoU gains and failures.

\paragraph{Implications for aerial foundation-model transfer.}
Overall, the results support three practical conclusions for transferring frozen visual foundation models to UAV 3D perception.  First, V-JEPA~2.1 features are useful for dense aerial geometry, but the transfer is not automatic and needs further investigation: a spatially structured decoder and metric metadata conditioning are needed.  Second, temporal/video features are valuable but must be converted into target-frame dense representations carefully; simple tubelet folding or local attention is not enough to guarantee consistent gains across altitude and scene.  Third, SSC performance depends on the compatibility between the geometry prior and the lifting/fusion architecture.  A strong MDE model can improve SSC, but only when its depth distribution and context features are aligned with the 3D view transformation and semantic decoder.

\subsection{Limitations}
\label{sec:limitations}

The study has several limitations.  First, the V-JEPA~2.1 backbone is frozen throughout.  This is useful for isolating representation transfer, but it does not test whether careful end-to-end finetuning could improve depth, temporal aggregation, or SSC.  Second, the depth heads operate on a $24\times24$ V-JEPA feature grid before dense decoding.  Fine structures in aerial scenes may require higher-resolution features or stronger feature pyramids than those tested in Table~\ref{tab:depth_block_flows}.  Third, the video experiments use a short four-frame clip with unknown frequency at which the images were taken.   and the DTA adapter defined in Eq.~\eqref{eq:dta}.  This is a controlled probe of spatio-temporal features, but it is not a complete exploration of temporal geometry, long-range motion, or learned frame selection.

Fourth, The poor performance of the unconditioned MVS SSC variant should therefore not be interpreted as a general failure of multiview stereo for aerial SSC.  Instead, it shows that cost-volume construction, feature resolution, altitude/intrinsics conditioning, and lifting-probability calibration are all critical.  Finally, the SSC results remain low in absolute semantic mIoU, even when they improve over the published OccuFly baselines in Table~\ref{tab:ssc_occufly_comparison}.

\subsection{Conclusion}
\label{sec:conclusion}

This paper studied whether frozen V-JEPA~2.1 image and video features transfer to aerial metric depth estimation and semantic scene completion on OccuFly.  The experiments show that frozen V-JEPA~2.1 representations contain useful geometric information, but that this information is not sufficient by itself.  Linear probing provides a meaningful baseline, while DPT-style dense decoding and FiLM conditioning with altitude and camera intrinsics are necessary for strong metric depth.  The best standard-split V-JEPA~2.1 depth result is obtained by FiLM\_DPT\_RGB in Table~\ref{tab:mde_comparison}, and the altitude-interpolation result in Table~\ref{tab:occufly-depth-30-50} further supports the value of physically meaningful metadata.

For semantic scene completion, the results show that depth priors improve aerial SSC.  The best overall SSC result is achieved by the FoundationSSC-style RGB-context variant with FiLM-conditioned depth, while the CGFormer comparison shows that replacing DAV2 depth with V-JEPA~2.1 FiLM depth can improve semantic mIoU with nearly unchanged occupied-space IoU.  The Pose-MVS experiments demonstrate that explicit multiview geometry is promising but sensitive: without altitude/intrinsics conditioning the MVS branch degrades severely, while FiLM-conditioned MVS recovers a substantial part of the gap.

The central conclusion is therefore that aerial 3D scene understanding requires more sophisticated downstream decoders to meaningfuly utilize the frozen V-JEPA~2.1 features. Frozen V-JEPA~2.1 provides a strong representation, but aerial metric reasoning benefits from decoder structure, calibrated metadata, and geometry-aware lifting.  Future work should focus on higher-resolution dense features, better target/source selection for UAV video and stronger cost-volume regularization.